\section{Inicio rapido}

Este manual explica como usar el RCE CQL en una clase. La idea es simple:
el docente levanta una URL, los alumnos entran sin login y cada navegador queda
asociado a su propio sandbox anonimo.

\subsection{Levantar el entorno local}

Desde Fedora, ubicarse en la carpeta del proyecto:

\begin{lstlisting}
cd ~/Desktop/RCE-CQL
cp .env.example .env
docker compose --profile local-hapi --profile local-translator up -d --build
\end{lstlisting}

Verificar servicios:

\begin{lstlisting}
docker ps
curl -s http://localhost:3000/api/v1/health/ready | jq
\end{lstlisting}

Abrir el RCE:

\begin{lstlisting}
http://localhost:5173
\end{lstlisting}

\rceScreenshot{01-docker-servicios.png}{Servicios Docker levantados para la demo local.}
\rceScreenshot{02-health-ready.png}{Health check del backend con HAPI y traductor disponibles.}

\subsection{Cargar pacientes sinteticos}

Si HAPI local esta vacio, cargar pacientes Synthea:

\begin{lstlisting}
docker compose --profile local-hapi --profile local-translator --profile seed-data run --rm synthea-seed
\end{lstlisting}

Confirmar que existen pacientes:

\begin{lstlisting}
curl -sG http://localhost:8080/fhir/Patient \
  --data-urlencode '_summary=count' | jq '.total'
\end{lstlisting}

\subsection{URLs importantes}

\begin{longtable}{>{\ttfamily}p{0.38\linewidth}p{0.52\linewidth}}
\toprule
URL & Uso \\
\midrule
http://localhost:5173 & Frontend del RCE. \\
http://localhost:3000/api/v1/health/ready & Estado del backend y dependencias. \\
http://localhost:3000/api/v1/cds-services & Discovery CDS Hooks. \\
http://localhost:8080/fhir & HAPI FHIR local. \\
http://localhost:8081 & CQL Translation Service local. \\
\bottomrule
\end{longtable}
